\documentclass[11pt]{article}
\usepackage{Style}

\nocite{*} % Comentar si quiero citar
%\addbibresource{bibliografia.bib} % Quitar el comentado si quiero usar bibliografia

\begin{document}

\begin{titlepage} % Portada con información necesaria
    \hspace*{-2cm}
    \raisebox{-2cm}[0pt][0pt]{
        \includegraphics[width=0.4\textwidth]{logo_universidad2.png}
    }

    \vspace{3.5cm}
    \begin{center}
        {\Large Pontificia Universidad Católica de Chile}
        
        \vspace{2cm}
        {\Huge \bfseries 
        Apuntes de Geometría Algebraica
        }
        
        \vspace{1.5cm}
        {\Large Benjamín Mateluna}
        
        \vfill
        {\large \today \par}
    \end{center}
\end{titlepage}

\tableofcontents % Índice de contenidos

\newpage
\section*{Introducción}
\addcontentsline{toc}{section}{Introducción}

\noindent La intención de estas notas o apuntes, es recopilar lo aprendido de geometría algebraica
durante la carrera y el magister; buscan servir como un banco de datos al cual le puedo consultar
cualquier duda que tenga sobre esta rama de las matemáticas. Además, lo utilizare como una manera
de estudiar el ramo que estoy cursando "Selected Topics on Complex Algebraic Varieties".

\vspace{1mm}
\noindent El propósito o enfoque inicial serán contenidos introductorios de geometría algebraica,
a medida que el curso avance, los temas irán evolucionando. Estos últimos también se pueden 
expandir en un futuro, a medida que vaya aprendiendo mas. Los apuntes están pensados para una
evolución constante.

\newpage
\section{Categorías y Funtores}
\emph{Esta sección se hará cuando tenga mas tiempo...}

\newpage
\section{Variedades Afines}
\noindent Empezaremos con la noción mas básica y sencilla de entender en la geometría algebraica, 
posiblemente la más cercana a la intuición geométrica. Comenzamos fijando un cuerpo $k$ 
algebraicamente cerrado, usualmente se piensa en $k=\C$.

\begin{dfn}
    El $n-$espacio afín sobre $k$, denotado por $\A_{k}^{n}$, es el conjunto de las $n-$tuplas de
    elemento de $k$, es decir, $k^{n}$. Un elemento $p=(a_{1},\cdots,a_{n})\in\A_{k}^{n}$ se dice 
    punto y $a_{i}$ coordenada de $p$.
\end{dfn}

\noindent El espacio afín, por ahora, lo pensamos únicamente como un conjunto y no con su 
estructura usual de espacio vectorial. Más adelante le daremos una estructura geométrica al 
$n-$espacio afín.

\begin{dfn}
    Sea $T\subseteq k[x_{1},\cdots,x_{n}]$, se define
    \begin{equation*}
        V(T):=\{p\in\A_{k}^{n}:f(p)=0,\forall f\in T\}
    \end{equation*}
    el lugar de ceros.
\end{dfn}

\noindent Cuando $T=\{f_{1},\cdots,f_{n}\}$ simplemente se escribe $V(f_{1},\cdots,f_{n})$.

\vspace{1mm}
\obs Claramente si $I=(T)$, entonces $V(I)=V(T)$. Más aún, como $k[x_{1},\cdots,x_{n}]$ es un 
anillo noetheriano, todo ideal es finitamente generado. Así, $V(T)$ se puede expresar como el 
lugar de ceros de un número finito de polinomios.

\begin{dfn}
    Un conjunto $V\subseteq\A_{k}^{n}$ se dice variedad algebraica afín si $V=V(T)$ para algún 
    $T\subseteq k[x_{1},\cdots,x_{n}]$.
\end{dfn}

\noindent En algunos textos, lo anterior se define simplemente como conjunto algebraico afín y
reservan el nombre de variedad algebraica afín a conjuntos algebraicos irreducibles, en 
nuestro caso, variedades algebraicas irreducibles, a saber

\begin{dfn}
    Una variedad algebraica afín $V$ se dice irreducible si no se puede expresar como la unión
    $V=V_{1}\cup V_{2}$ de dos subconjuntos propios y cada uno de ellos una variedad algebraica.
\end{dfn}

\noindent Pensando en teoría categorías, hemos definido nuestros objetos, faltan los morfismos.

\begin{dfn}
    Sean $V\subseteq\A_{k}^{n}$ y $W\subseteq\A_{k}^{n}$ variedades afines. Una función 
    $\varphi:V\to W$ es un morfismo de variedades afines si es la restricción de un mapa 
    polinomial.
\end{dfn}

\noindent Un mapa polinomial es aquel que en cada coordenada esta dado por polinomios.

\begin{ejm}
    Consideremos $V=V(x^{2}-y)$, el mapa $\varphi:\C\to V$ dado por
    \begin{equation*}
        \varphi(z)=(z,z^{2})
    \end{equation*}
    es un morfismo de variedades afines, de hecho, es un isomorfismo.
\end{ejm}

\noindent De la misma manera que relacionamos un conjunto de polinomios con un conjunto de 
puntos en el espacio afín, podemos hacer el proceso en reversa

\begin{dfn}
    Dado $V\subseteq\A_{k}^{n}$, definimos el ideal de $V$ en $k[x_{1},\cdots,x_{n}]$, como
    \begin{equation*}
        I(V):=\{f\in k[x_{1},\cdots,x_{n}]:f|_{V}\equiv0\}
    \end{equation*}
\end{dfn}

\noindent Es facil ver que $I(V)$ es un ideal, más aún, si $f^{m}\in I(V)$ entonces $f\in I(V)$, 
en otras palabras, el ideal de un conjunto es radical.

\begin{ejm}
    Consideremos $V=V(x^{2})=\{0\}\subseteq\A_{k}^{1}$. Notemos que $I(V)=(x)$.
\end{ejm}

\noindent Surge entonces la pregunta una pregunta natural ¿Cuál es exactamente la correspondencia
entre conjuntos de polinomios y conjuntos en el espacio afín? Como veremos mas adelante, el 
Nullstellensatz de Hilbert nos dará una respuesta precisa, por ahora, veamos el siguiente 
resultado

\newpage
\begin{prop}
    Una variedad afín $V\subseteq\A_{k}^{n}$ es irreducible si y solo si $I(V)$ es un ideal primo.
\end{prop}

\begin{proof}
    Supongamos que $I(V)$ no es primo, entonces existen $f_{1},f_{2}$ polinomios tales que 
    $f_{1}f_{2}\in I(V)$, pero $f_{i}\not\in I(V)$. Notar que
    \begin{equation*}
        V=(V\cap V(f_{1}))\cup(V\cap V(f_{2}))
    \end{equation*}
    pero $(V\cap V(f_{i}))\neq V$, ya que existe $p_{i}\in V$ tal que $f_{i}(p_{i})\neq0$. Luego, 
    $V$ es reducible.

    \vspace{1mm}
    Supongamos que $V$ es reducible, entonces
    \begin{equation*}
        V=V_{1}\cup V_{2}
    \end{equation*}
    con $V_{i}\neq V$ variedades afines. Entonces existe un polinomio $f_{i}$ tal que $f_{i}(p)=0$ 
    para todo $p\in V_{i}$, pero no para todo punto en $V$. Notemos que $f_{1}f_{2}\in I(V)$, pero 
    $f_{i}\not\in I(V)$. Concluimos que $I(V)$ no es primo.
\end{proof}

%\printbibliography % Quitar el comentado si quiero usar bibliografia

\end{document}
